 \section{Introduction}

Natural language inference (NLI; also known as recognizing textual entailment, or RTE)
is a widely-studied task in natural language processing, to which many complex semantic tasks, such as question answering and text summarization, can be reduced \cite{Dagan:2006}.\resolved{\nascomment{can someone convince me this is true?  I don't see how QA or summarization reduces to classifying sentence pairs.}.}
Given a pair of sentences, a premise $p$ and a hypothesis $h$, the goal is to determine whether or not $p$ semantically entails $h$.
% Given a pair of sentences---a premise $p$ and a hypothesis $h$---the goal is to determine whether $p$ entails $h$; i.e., given that $p$ is true, is $h$ also true? % TODO: This sentence is a bit odd (since it ignores the neutral/contradiction distinction), and is partially redundant with the next paragraph. Cut/rework?

The problem of acquiring large amounts of labeled inference data was addressed by \citet{Bowman:2015}, who devised a method for crowdsourcing high-agreement entailment annotations en masse, creating the SNLI and later the genre-diverse MultiNLI \cite{Williams:2017} datasets.
In this process, crowd workers are presented with a premise $p$ drawn from some corpus (e.g., image captions), and are required to generate three new sentences (hypotheses) based on $p$, according to one of the following criteria:
% 
% \noindent
% ~~~\textbf{Entailment}: $h$ is definitely true given $p$ 

% \noindent
% ~~~\textbf{Neutral}: $h$ might be true given $p$

% \noindent
% ~~~\textbf{Contradiction}: $h$ is definitely NOT true given $p$
% 
% \noindent
% \com{
\begin{table}[h]
\vspace{-0.5em}
\begin{tabular}{ll}
%\multicolumn{2}{c}{} \\
% \hline
\textbf{Entailment}    & $h$ is definitely true given $p$ \\
\textbf{Neutral}       & $h$ might be true given $p$ \\
\textbf{Contradiction} & $h$ is definitely {\bf not} true given $p$ \\
% \hline
%\multicolumn{2}{c}{} \\
\end{tabular}
\vspace{-1em}
\end{table}
% }
% The simple nature of these instructions allowed the annotation process to scale, resulting in the collection of roughly one million high-agreement examples (570K in SNLI and 433K in MultiNLI).\footnote{Both datasets were partially re-annotated in a validation phase that ensured the quality of the generated hypotheses.\sw{what is this footnote really trying to say?}}
% The release of SNLI in 2015 cleared the way for new deep learning models \cite[\textit{inter alia}]{Parikh:2016,Chen:2017,Gong:2017} that quickly outperformed the simpler pre-neural models \cite{Bowman:2015}.
\noindent
In this paper, we observe that hypotheses generated by this crowdsourcing process contain artifacts that can help a classifier detect the correct class \textit{without} ever observing the premise (\secref{artifacts}).
% We find that by reading the hypothesis $h$ alone 
% , text categorization models can achieve over 67\% and 55\% accuracy on SNLI and MultiNLI, respectively, 
% compared to a 35\% most-frequent-class baseline (\secref{artifacts}). 
% This result implies that these annotation artifacts indeed emerge from the crowdsourcing protocol, and that they are abundant in NLI datasets.
% without learning general entailment phenomena.
% \sam{I don't think you make the case that these artifacts are unrelated to NLI, but there is probably something non-trivial you can say here... rework?}

\begin{table*}[t]
\centering
{
\begin{tabular}{ll} \toprule
\textbf{Premise} & A woman selling bamboo sticks talking to two men on a loading dock. \\\midrule
\textbf{Entailment} & There are {\bf at least} three {\bf people} on a loading dock. \\
\textbf{Neutral} & A woman is selling bamboo sticks {\bf to help provide for her family.} \\
\textbf{Contradiction} & A woman is {\bf not} taking money for any of her sticks. \\\bottomrule
\end{tabular}
}
\caption{An instance from SNLI that illustrates the artifacts that arise from the annotation protocol. 
A common strategy for generating entailed hypotheses is to remove gender or number information. 
Neutral hypotheses are often constructed by adding a purpose clause. 
Negations are often introduced to generate contradictions.}
\label{tab:example}
\end{table*}

A closer look suggests that the observed artifacts are a product of specific annotation strategies and heuristics that crowd workers adopt. %, perhaps to improve their efficiency.
We find, for example, that entailed hypotheses tend to contain gender-neutral references to people, purpose clauses are a sign of neutral hypotheses, and negation is correlated with contradiction (\secref{analysis}). 
Table~\ref{tab:example} shows a single set of instances from SNLI that demonstrates all three phenomena.

We re-evaluate high-performing NLI models on the subset of examples on which our hypothesis-only classifier failed, which we consider to be ``hard'' (\secref{sota}). 
Our results show that the performance of these models on the ``hard'' subset is dramatically lower than their performance on the rest of the instances. 
This suggests that, despite recently reported progress, natural language inference remains an open problem.
